% Springer-compatible LaTeX template for APJM / MIR submission
% Compiled from Markdown via pandoc.
% Document class: `article` (Springer Nature's sn-jnl is not bundled; this template
% follows Springer's general manuscript guidelines and is compatible with both
% Springer Nature LaTeX template and Word-track submission portals).

\documentclass[12pt,a4paper]{article}

% --- Geometry & spacing ---
\usepackage[margin=2.5cm]{geometry}
\usepackage{setspace}
\onehalfspacing

% --- Encoding & fonts ---
% Conditional setup: under pdflatex use legacy inputenc; under xelatex/lualatex
% use fontspec for native Unicode support (needed for β, −, × symbols).
\usepackage{iftex}
\ifPDFTeX
  \usepackage[utf8]{inputenc}
  \usepackage[T1]{fontenc}
  \usepackage{lmodern}
\else
  \usepackage{fontspec}
  \setmainfont{Liberation Serif}[
    Scale=1.0,
    BoldFont={Liberation Serif Bold},
    ItalicFont={Liberation Serif Italic},
    BoldItalicFont={Liberation Serif Bold Italic}
  ]
  % Liberation Mono carries the full Vietnamese diacritic set; DejaVu Sans
  % Mono does not (renders Ấ/ầ/ế/ố as missing characters in code blocks).
  \setmonofont{Liberation Mono}[Scale=0.85]
  % Map common Unicode math/typographic symbols that may appear in markdown
  % body text. Liberation Serif lacks them, so route via amssymb/math mode.
  % (We deliberately avoid unicode-math, which conflicts with the amssymb
  % package that pandoc auto-loads in the preamble.)
  \usepackage{newunicodechar}
  \newunicodechar{∈}{\ensuremath{\in}}
  \newunicodechar{∉}{\ensuremath{\notin}}
  \newunicodechar{≤}{\ensuremath{\leq}}
  \newunicodechar{≥}{\ensuremath{\geq}}
  \newunicodechar{≠}{\ensuremath{\neq}}
  \newunicodechar{≈}{\ensuremath{\approx}}
  \newunicodechar{∞}{\ensuremath{\infty}}
  \newunicodechar{∑}{\ensuremath{\sum}}
  \newunicodechar{∏}{\ensuremath{\prod}}
  \newunicodechar{∗}{\ensuremath{\ast}}
  \newunicodechar{·}{\ensuremath{\cdot}}
  \newunicodechar{−}{\ensuremath{-}}
  \newunicodechar{×}{\ensuremath{\times}}
  \newunicodechar{÷}{\ensuremath{\div}}
  \newunicodechar{±}{\ensuremath{\pm}}
  % Typographic checkmarks/crosses (used as bullet markers in markdown).
  \newunicodechar{✓}{\ensuremath{\checkmark}}
  \newunicodechar{✗}{\ensuremath{\times}}
  \newunicodechar{⚠}{\textbf{!}}
  \newunicodechar{→}{\ensuremath{\rightarrow}}
  \newunicodechar{←}{\ensuremath{\leftarrow}}
  \newunicodechar{⇒}{\ensuremath{\Rightarrow}}
  % Allow URL line breaks at slashes (fix Overfull \hbox in long URLs)
  \usepackage{xurl}
\fi

% --- Math ---
\usepackage{amsmath,amssymb,amsfonts,mathtools}
\usepackage{bm}

% --- Tables ---
\usepackage{booktabs}
\usepackage{longtable}
\usepackage{multirow}
\usepackage{array}
\usepackage{tabularx}

% --- Figures ---
\usepackage{graphicx}
\usepackage{float}
\usepackage{caption}
\captionsetup{font=small,labelfont=bf}

% --- Misc text ---
\usepackage{microtype}
\usepackage{enumitem}
\setlist{nosep,leftmargin=1.5em}
\usepackage{textcomp}
\usepackage{xcolor}
\usepackage{ulem}\normalem

% --- Code block support (pandoc Shaded / Highlighting environments) ---
% Required when source markdown contains fenced code blocks (```...```).
\usepackage{fancyvrb}
\usepackage{framed}
\definecolor{shadecolor}{RGB}{248,248,248}
\newenvironment{Shaded}{\begin{snugshade}}{\end{snugshade}}
% Highlighting environment for syntax-highlighted code
\providecommand{\BuiltInTok}[1]{#1}
\providecommand{\CommentTok}[1]{\textit{#1}}
\providecommand{\KeywordTok}[1]{\textbf{#1}}
\providecommand{\NormalTok}[1]{#1}
\providecommand{\StringTok}[1]{\textit{#1}}
\providecommand{\OperatorTok}[1]{#1}
\providecommand{\DataTypeTok}[1]{#1}
\providecommand{\FunctionTok}[1]{#1}

% --- Links & references ---
\usepackage[hidelinks,breaklinks=true]{hyperref}
\usepackage{cleveref}

% --- Bibliography (APA7 via natbib + apalike) ---
% APJM/MIR accept APA7. natbib's apalike close approximation; for full APA7
% compliance use biblatex with apa style.
\usepackage[round,sort]{natbib}
\bibliographystyle{apalike}

% --- Pandoc compatibility shims ---
\providecommand{\tightlist}{\setlength{\itemsep}{0pt}\setlength{\parskip}{0pt}}
\providecommand{\passthrough}[1]{\lstset{mathescape=false}#1\lstset{mathescape=true}}

% --- Title page (filled in per paper) ---
\title{p7\_capstone submission to JIBS}
\author{Author details withheld for blind review}
\date{2026-05-21}

\begin{document}

\maketitle

\hypertarget{internationalization-and-firm-performance-across-49-asian-and-pacific-economies-institutional-regimes-digital-capabilities-and-managerial-characteristics-as-contingency-factors}{%
\section{Internationalization and Firm Performance across 49 Asian and
Pacific Economies: Institutional Regimes, Digital Capabilities, and
Managerial Characteristics as Contingency
Factors}\label{internationalization-and-firm-performance-across-49-asian-and-pacific-economies-institutional-regimes-digital-capabilities-and-managerial-characteristics-as-contingency-factors}}

\textbf{Version:} Working manuscript (May 2026)\\
\textbf{Target journal:} Journal of International Business Studies
(JIBS)\\
\textbf{Status:} Empirical results complete; under revision for
submission

\begin{center}\rule{0.5\linewidth}{0.5pt}\end{center}

\hypertarget{abstract}{%
\subsection{Abstract}\label{abstract}}

\textbf{Purpose.} This study examines whether the
internationalisation-performance (I-P) relationship holds across a large
and institutionally diverse set of Asian and Pacific economies and
identifies the conditions under which firms convert export intensity
into productive performance.

\textbf{Design/methodology/approach.} Using microdata from 91,982 firms
across 49 economies and 102 country-year waves of the World Bank
Enterprise Survey (2003-2025), spanning all six Institutional Capability
and Resource Vulnerability (ICRV) regime groups, we estimate a
hierarchical sequence of ordinary least squares (OLS) models with HC1
robust standard errors and three-way interaction specifications.

\textbf{Findings.} Three results emerge. First, the I-P relationship is
robustly inverted-U with a turning point at approximately 36\% of
foreign sales to total sales (Lind-Mehlum test p \textless{} .001;
confirmed in the full three-way moderation model at turning point
34.6\%, p = .002). Second, digital adoption (DAI) both raises the
performance level and significantly reshapes the I-P curve, generating
compressed ascending-limb returns at low export intensity and attenuated
performance decline at high export intensity. Third, institutional
regime (ICRV) systematically moderates both slope and curvature, with
stronger per-unit digital returns in weaker institutional environments.
Technological capability raises the performance level without reliably
reshaping the curve; managerial experience and female top management
improve performance.

\textbf{Originality/value.} The study offers a parsimonious resolution
to the long-standing heterogeneity in Asian I-P findings: the
relationship is universally inverted-U, but regime quality and digital
capability jointly determine where the turning point falls. The
49-economy pool is the largest WBES-based I-P test in Asia and the
Pacific to date, establishing a benchmark against which country-specific
results (Vietnam, Singapore, China) can be calibrated.

\textbf{Keywords:} internationalisation-performance; inverted-U;
institutional regime; digital capability; Asia-Pacific; World Bank
Enterprise Surveys.

\textbf{JEL Classification:} F23; O33; D22; L25; O53.

\textbf{Paper type:} Research paper.

\begin{center}\rule{0.5\linewidth}{0.5pt}\end{center}

\hypertarget{1-introduction}{%
\subsection{1. Introduction}\label{1-introduction}}

A central question in international business research concerns whether
and how export internationalization translates into firm performance (Lu
\& Beamish, 2004; Johanson \& Vahlne, 2009). Meta-analyses document a
mean correlation of approximately r = .07 across several thousand
studies (Bausch \& Krist, 2007; Kirca et al., 2012; Marano et al.,
2016), yet the variance across studies vastly exceeds what sampling
error can explain. In Asia, a region that accounts for over half of
global manufacturing exports (UNCTAD, 2023), the evidence is especially
fragmented: prior studies report positive relationships (Pangarkar,
2008; Korea), inverted-U curves (Chiao et al., 2006; Taiwan), or null
and negative effects (Mondal et al., 2022; India). The standard
explanation invokes context as a moderator (Marano et al., 2016; Kirca
et al., 2012), but empirical tests simultaneously across institutionally
diverse Asian economies remain rare.

This paper addresses that gap. We pool 91,982 firm-observations from 49
Asian and Pacific economies spanning 2003--2025, harmonised from the
World Bank Enterprise Survey (WBES). Our sample spans all six ICRV
regime groups, from advanced innovation-driven economies (Singapore,
Korea, Taiwan, HongKong, Israel, Cyprus) to Pacific Small Island
Developing States (Vanuatu, Solomon Islands, Fiji, Tonga, Samoa,
Kiribati, Papua New Guinea, Timor-Leste) and Gulf advanced-resource
economies (Bahrain, Kuwait, Qatar, Brunei, Saudi Arabia, Oman),
classified through the Institutional Capability and Resource
Vulnerability (ICRV) framework that anchors the broader dissertation
programme of which this study is part.

Our contributions are threefold. First, we provide the largest
WBES-based test of the inverted-U I--P hypothesis in Asia, confirming a
turning point at approximately 36\% foreign sales intensity across all
specifications (N = 84,910--29,840), robust even in the full three-way
moderation model (M11: TP=34.6\%, LM p=.002). Second, we document that
digital adoption (DAI) is the primary capability moderator that reshapes
the I--P curve, generating significant curve compression effects
(FSTS×DAI p\textless.001) that emerge even without controlling for
manager characteristics. Technological capability (TCI) raises the
performance level but does not reliably reshape the ascending limb
across this sample. Third, we document that institutional regime (ICRV
group) systematically moderates both the slope and curvature of the I--P
function, with DAI×ICRV interactions (p=.012) indicating that digital
capabilities deliver stronger per-unit performance returns in weaker
institutional environments.

The remainder of the paper is organised as follows. Section 2 develops
the theoretical framework and formal hypotheses. Section 3 describes
data and methods. Section 4 presents results. Section 5 discusses
implications and limitations. Section 6 concludes.

\begin{center}\rule{0.5\linewidth}{0.5pt}\end{center}

\hypertarget{2-theoretical-framework-and-hypotheses}{%
\subsection{2. Theoretical Framework and
Hypotheses}\label{2-theoretical-framework-and-hypotheses}}

\hypertarget{21-theoretical-foundations}{%
\subsubsection{2.1 Theoretical
Foundations}\label{21-theoretical-foundations}}

We draw on four theoretical traditions. The Uppsala model (Johanson \&
Vahlne, 1977, 2009) describes internationalization as a sequential
process in which firms accumulate experiential knowledge, progressively
reducing psychic distance and coordination costs. The Resource-Based
View (Barney, 1991; Wernerfelt, 1984) positions firm-level capabilities,
technological, digital, and managerial, as the mediating mechanism
linking export exposure to performance outcomes. Institutional Theory
(North, 1990; Scott, 2008) directs attention to how country-level formal
and informal institutions shape the cost structure of cross-border
transactions. Upper Echelons Theory (Hambrick \& Mason, 1984; Hambrick,
2007) emphasises that top management team attributes moderate how firms
respond to strategic opportunities and constraints.

\hypertarget{22-the-non-linear-ip-relationship-h1}{%
\subsubsection{2.2 The Non-Linear I--P Relationship
(H1)}\label{22-the-non-linear-ip-relationship-h1}}

Contractor et al. (2003) formalised the three-stage S-curve hypothesis:
early internationalization carries fixed entry costs that depress
performance; intermediate internationalization yields learning and scale
economies; high internationalization may again impose coordination and
agency costs. Empirically, however, the data more consistently support
the inverted-U, a special case in which the third stage does not emerge
at the export intensities typically observed in developing-economy WBES
data (Gomes \& Ramaswamy, 1999; Lu \& Beamish, 2004).

In the Asian context, three country-specific studies in our broader
dissertation programme provide prior evidence. For Vietnam (P3), the
inverted-U is confirmed with a turning point at 39--46\% using a
WBES-based instrumental variable design. For China (P5), an inverted-U
turning point of 47--49\% is replicated across specifications with
Paternoster cross-cohort tests (p = .545). For Singapore (P4), the
relationship is predominantly positive with a very late turning point
(\textasciitilde88.6\%), consistent with near-monotonic returns in a
high-institutional, digitally saturated economy. Pooled across a much
larger and institutionally heterogeneous sample, we expect the
inverted-U to emerge at an intermediate turning point.

\begin{quote}
\textbf{H1:} The relationship between export intensity (FSTS) and firm
labour productivity (ln LP) is an inverted-U, with a statistically
confirmed turning point at intermediate export intensities, in a pooled
sample of Asian and Pacific firms.
\end{quote}

\hypertarget{23-technological-capability-as-a-curve-amplifier-h2}{%
\subsubsection{2.3 Technological Capability as a Curve Amplifier
(H2)}\label{23-technological-capability-as-a-curve-amplifier-h2}}

Technological capability, operationalised as a composite of quality
certification, foreign-technology adoption, process innovation, and
R\&D, raises absorptive capacity (Cohen \& Levinthal, 1990), enabling
firms to extract greater performance value per unit of export exposure.
Critically, we expect TCI to operate as a \emph{curve amplifier}: at low
export intensities, high-TCI firms gain more per unit of
internationalization (steeper ascending limb); at high intensities, they
experience sharper diminishing returns as coordination demands exceed
even their superior capabilities (steeper descending limb). The net
effect is a shift and steepening of the inverted-U, not just a parallel
upward shift.

This prediction is consistent with P3 Vietnam (β(TCI) = +0.179,
IV-identified) and P5 China (β(TCI) increasing from +0.28 to +0.43
across cohorts), both showing TCI as a positive level shifter. P4
Singapore shows negligible TCI moderation of the curve, likely because
near-universal adoption in a high-ICRV economy leaves insufficient
variance for moderation.

\begin{quote}
\textbf{H2:} TCI amplifies the I--P curve, raising the average
performance level and steepening both the ascending and descending limbs
of the inverted-U.
\end{quote}

\hypertarget{24-digital-adoption-as-a-level-enhancer-h3}{%
\subsubsection{2.4 Digital Adoption as a Level Enhancer
(H3)}\label{24-digital-adoption-as-a-level-enhancer-h3}}

Digital adoption, measured through website presence and
electronic-payment share, has been associated with reduced communication
and transaction costs (Verhoef et al., 2021). In the CDCM
(Context-Contingent Digital Capability Model) developed in the
dissertation, DAI\textquotesingle s effect depends on the interaction of
regime quality, FSTS level, and digital market saturation (Đỗ \& Phan,
2024). Specifically, in institutionally advanced economies where digital
infrastructure is mature, DAI provides universal performance lifts (as
in P4 Singapore: β(DAI) positive and significant). In emerging-economy
contexts with substantial selection into digital adoption (P3 Vietnam:
IV-estimated β(DAI) ≈ 0, consistent with selection bias in OLS),
DAI\textquotesingle s curve-shaping role is attenuated. Across the
pooled multi-economy sample, we expect DAI to raise performance levels
broadly across the sample (level effect). In a fully specified model
that also accounts for managerial quality, which correlates with digital
adoption, the residual DAI variation may reveal a
capability-complementarity curve effect: digitally equipped firms at low
FSTS face sunk adoption costs that compress near-term productivity,
while digital infrastructure supports coordination efficiency at high
FSTS where cross-border management demands are greatest.

\begin{quote}
\textbf{H3:} DAI has a positive main effect on labour productivity and,
when managerial characteristics are controlled, reshapes the I--P curve
by compressing the ascending limb at low export intensities and
attenuating performance decline at high export intensities.
\end{quote}

\hypertarget{25-managerial-characteristics-h4}{%
\subsubsection{2.5 Managerial Characteristics
(H4)}\label{25-managerial-characteristics-h4}}

Upper Echelons Theory predicts that top manager experience and gender
composition affect strategic cognition and decision-making (Hambrick \&
Mason, 1984). Manager experience reduces coordination errors in
cross-border operations; female top management is associated with
stakeholder orientation, relationship capital, and risk management, all
relevant in institutionally heterogeneous export markets (Richard et
al., 2019). We expect both to raise the performance level for any given
FSTS. However, given the breadth of the multi-country sample, firm-level
managerial differences are unlikely to systematically reshape the
aggregate FSTS--performance curve, whose shape is primarily determined
by economy-wide institutional and market-maturity factors.

\begin{quote}
\textbf{H4:} Top manager experience and female top manager status have
positive main effects on firm performance, but do not significantly
moderate the FSTS--performance curve shape.
\end{quote}

\hypertarget{26-institutional-regime-as-a-turning-point-shifter-h5}{%
\subsubsection{2.6 Institutional Regime as a Turning-Point Shifter
(H5)}\label{26-institutional-regime-as-a-turning-point-shifter-h5}}

North (1990) established that formal institutions reduce transaction
costs; stronger institutions lower the cost of cross-border coordination
and thus shift the performance-maximising export intensity (the
inverted-U turning point). In higher-quality institutional environments,
firms need a lower export share to recoup internationalization costs, so
the turning point occurs earlier. Equivalently, holding FSTS constant,
firms in stronger institutional regimes earn higher returns per unit of
internationalization.

We operationalise institutional regime through the ICRV classification
system developed in the dissertation: six groups ranging from Group I
(Advanced innovation-driven: Singapore, Korea, Taiwan, Israel, Cyprus)
to Group VI (SIDS and small open economies: Vanuatu, Solomon Islands,
Timor-Leste). Higher ICRV group numbers denote weaker institutions. We
expect positive FSTS×ICRV and negative FSTS²×ICRV interactions in Group
III--VI compared to Group I (stronger institutions = higher Group I
coefficient), which is observationally equivalent to the turning point
occurring later in lower-quality institutional environments.

\begin{quote}
\textbf{H5:} The ICRV institutional regime positively moderates the I--P
relationship: firms in higher-quality institutional environments (lower
ICRV group number) achieve the performance-maximising export intensity
at lower FSTS levels.
\end{quote}

\hypertarget{27-regime-contingent-digital-effects-h6}{%
\subsubsection{2.7 Regime-Contingent Digital Effects
(H6)}\label{27-regime-contingent-digital-effects-h6}}

The Context-Contingent Digital Capability Model (CDCM) developed in this
dissertation programme predicts that the I--P curve-reshaping role of
digital adoption is not uniform across institutional environments. In
institutionally advanced economies (ICRV Group I--II), digital
infrastructure is near-universal and close to the competitive baseline,
the marginal performance differentiation from website and e-payment
adoption is accordingly lower. In institutionally weaker economies (ICRV
Group IV--VI), digital adoption remains differentiating: digital market
access tools lower the liability-of-foreignness cost and provide
cross-border coordination support that institutional voids cannot
substitute. The DAI×ICRV interaction therefore tests whether digital
capabilities\textquotesingle{} curve-reshaping potency (established in
H3) varies systematically across the institutional spectrum,
specifically, whether compressing the ascending limb and buffering
performance decline at high FSTS are more pronounced in
weaker-institution economies.

This prediction is consistent with the digital-complementarity strand of
institutional theory (North, 1990; Stallkamp \& Schotter, 2021): where
formal institutions provide reliable contract enforcement, IP
protection, and market information, digital infrastructure adds
incrementally; where institutional voids prevail, digital tools provide
a non-trivial substitute for missing formal mechanisms. H6 therefore
does not predict stronger \emph{level} effects in weak-institution
contexts (which could reflect many confounds), but specifically predicts
a stronger \emph{curve-shaping} interaction between digital adoption and
export intensity, a conditional complementarity claim.

\begin{quote}
\textbf{H6:} The ICRV institutional regime moderates the
digital-adoption curve-reshaping effect: digital
adoption\textquotesingle s compressive effect on the I--P ascending limb
and its buffering effect at high FSTS (H3) are stronger in weaker
institutional environments (higher ICRV group number), where digital
infrastructure provides greater competitive differentiation in
cross-border market access.
\end{quote}

\begin{center}\rule{0.5\linewidth}{0.5pt}\end{center}

\hypertarget{3-data-and-methods}{%
\subsection{3. Data and Methods}\label{3-data-and-methods}}

\hypertarget{31-data-source}{%
\subsubsection{3.1 Data Source}\label{31-data-source}}

Data derive from the World Bank Enterprise Surveys (WBES), a stratified
random-sample survey of formal private-sector enterprises conducted
periodically across over 130 countries. We use microdata from 49 Asian
and Pacific economies across 102 country-year waves spanning 2003--2025.
This coverage encompasses six UN M.49 Asian sub-regions: East Asia
(China, HongKong, Korea, Mongolia, Taiwan), Southeast Asia (Brunei,
Cambodia, Indonesia, Laos, Malaysia, Myanmar, Philippines, Singapore,
Thailand, Timor-Leste, Vietnam), South Asia (Afghanistan, Bangladesh,
Bhutan, India, Maldives, Nepal, Pakistan, Sri Lanka), Central Asia
(Kazakhstan, Kyrgyz Republic, Tajikistan, Turkmenistan, Uzbekistan),
West Asia (Armenia, Bahrain, Cyprus, Iraq, Israel, Jordan, Kuwait,
Lebanon, Oman, Qatar, Saudi Arabia, Yemen), and Pacific/SIDS (Fiji,
Kiribati, Papua New Guinea, Samoa, Solomon Islands, Tonga, Vanuatu).

The final analytical dataset contains 91,982 firm-observations (before
applying outcome and FSTS non-missing requirements).
Oman\textquotesingle s available data (2003 vintage) predates the
digital-capability variables used in M6--M11; it contributes to M2 but
drops from M3+ due to missing controls. WBES designs are cross-sectional
within each wave, with standardised instruments ensuring common variable
definitions across economies and years. We apply three harmonisation
steps: (1) variable-priority mapping for candidate items across PICS3
(2009--2013), Standardised (2014--2018), and BREADY/BEE (2019--2025)
questionnaire generations; (2) encoding-robust reading (primary UTF-8
with Latin-1/CP1252 fallback); (3) deduplication, selecting the largest
file per (country, year) pair when multiple uploads exist.

\hypertarget{32-variables}{%
\subsubsection{3.2 Variables}\label{32-variables}}

\textbf{Dependent variable.} Labour productivity is operationalised as
log annual sales per permanent worker: ln(LP) = ln(sales / l1\_workers).
Sales are taken from WBES item n3 (or d2 where n3 is missing). Currency
effects are absorbed by country-year fixed effects in M5; without FE,
productivity levels vary in nominal terms across economies and years.

\textbf{Foreign sales intensity (FSTS).} Defined as the share of sales
from direct exports (d3c / 100), ranging from 0 to 1. We mean-centre
FSTS at the sample mean for quadratic and interaction models (fsts\_c);
the quadratic term fsts\_c² is computed from the centred value.

\textbf{Technological capability index (TCI\_z).} A two-item composite
(available in \textgreater95\% of files) of quality certification (b8:
has ISO or other certification) and foreign-technology adoption (e6:
uses foreign-licensed technology), standardised to mean 0, SD 1. The
z-score preserves cross-economy variation.

\textbf{Digital adoption index (DAI\_z).} Website presence (c22b / e1)
averaged with e-payment share (k33 / 100 where available), standardised
to mean 0, SD 1. In waves without e-payment items, DAI equals website
presence only.

\textbf{Manager characteristics.} Top manager experience in sector (b7,
in years) and a binary indicator for female top manager (b7a / b6a).

\textbf{Controls.} Female ownership share (b4 \textgreater{} 0.5:
majority female-owned, binary), foreign ownership percentage (b6a /
100), firm age (in years, current year − b5), and log permanent
employees (ln l1\_workers) as a size proxy.

\textbf{Institutional regime (ICRV group).} Integer 1--6 based on the
ICRV classification; used as a continuous moderator in M10/M11 and as a
group-level moderator for marginal effects plots.

\hypertarget{33-estimation-strategy}{%
\subsubsection{3.3 Estimation Strategy}\label{33-estimation-strategy}}

We estimate a hierarchical sequence of 11 models (Table 2). All models
use OLS with HC1 heteroskedasticity-robust standard errors (Long \&
Ervin, 2000), clustering by country-year. Country-year fixed effects
(M5) provide the most conservative benchmark. We prefer HC1 over
clustered SE given that WBES sampling is stratified within, not across,
country-years.

The inverted-U shape is formally confirmed using the Lind--Mehlum (2010)
test, which requires: (a) β₁ \textgreater{} 0, β₂ \textless{} 0, and (b)
the turning point lies strictly within the observed FSTS range, and (c)
the slope at the extremes satisfies the directional prediction. We
report the turning point as: TP = −β₁ / (2β₂) + mean(FSTS), where β₁ and
β₂ are estimated on the mean-centred FSTS\_c scale and the result is
expressed as a proportion (0--1) of total sales.

For moderation tests (M7--M11), we include interaction terms
FSTS×moderator and FSTS²×moderator. The joint significance of the two
interaction terms is tested with an F-test; we report both the joint
p-value and individual coefficient p-values.

\hypertarget{34-sample-sizes}{%
\subsubsection{3.4 Sample Sizes}\label{34-sample-sizes}}

The analytic sample for M2 (outcome + FSTS non-missing) is N = 84,910.
Adding controls reduces the sample to N = 38,342 (M3); country-year FE
model M5 uses the same sample. TCI-moderation models (M6, M7) use N ≈
38,051; adding DAI (M8) N ≈ 37,940; manager models (M9) N ≈ 35,568;
ICRV-moderation (M10) N ≈ 31,928; full three-way (M11) N ≈ 29,840. The
drop from M2 to M3/M5 reflects missingness in control variables
(primarily foreign ownership and firm age), which is more prevalent in
smaller economies and early WBES waves.

\begin{center}\rule{0.5\linewidth}{0.5pt}\end{center}

\hypertarget{4-results}{%
\subsection{4. Results}\label{4-results}}

\hypertarget{41-descriptive-statistics-and-correlations}{%
\subsubsection{4.1 Descriptive Statistics and
Correlations}\label{41-descriptive-statistics-and-correlations}}

The full analytic sample spans 49 economies and 102 country-year waves.
The mean FSTS across firms is approximately 12\% (median ≈ 0),
reflecting the fact that most WBES firms are domestically oriented;
exporters (FSTS \textgreater{} 0) represent approximately 18\% of the
sample. Labour productivity ranges from ln(LP) ≈ 10 (Tajikistan,
Timor-Leste) to ln(LP) ≈ 20 (Vietnam, Indonesia, reflecting
PPP-unadjusted nominal sales differences). Pairwise correlations among
focal variables are modest (\textbar r\textbar{} \textless{} .15 for all
pairs), and variance inflation factors in the full model are below 3.5,
ruling out multicollinearity concerns.

\hypertarget{42-main-effect-inverted-u-h1}{%
\subsubsection{4.2 Main Effect: Inverted-U
(H1)}\label{42-main-effect-inverted-u-h1}}

\emph{Table 2, Models M0--M5.}

The baseline linear model (M0) confirms that FSTS alone adds no
explanatory power (adjR² = .000). Adding the quadratic term FSTS² (M2)
yields β₁ = +1.316 (p \textless{} .001) and β₂ = −1.810 (p \textless{}
.001), consistent with an inverted-U. Following Haans, Pieters, and He
(2016), all four formal conditions for a genuine inverted-U are
satisfied in M2: (C1) β₁ = +1.316 (p \textless{} .001), confirming a
positive ascending limb; (C2) β₂ = −1.810 (p \textless{} .001),
confirming concavity; (C3) the turning point TP* ≈ 36.4\% lies well
within the observed FSTS range {[}0\%, 100\%{]}; and (C4) the marginal
effect of internationalization is positive at the lower data boundary
(FSTS ≈ 0\%, slope ≈ +1.32) and negative at the upper boundary (FSTS ≈
100\%, slope ≈ −2.30), with opposite signs confirmed by the Lind--Mehlum
(2010) u-test (joint p \textless{} .001). The inverted-U remains
significant when controls are added (M3: TP = 33.8\%, LM p \textless{}
.001) and when country-year fixed effects absorb all unobserved
time-varying country heterogeneity (M5: TP = 40.0\%, LM p \textless{}
.001, adjR² = .677). The consistency of the turning point across M2, M3,
and M5, ranging from 33.8\% to 40.0\%, provides strong evidence that the
inverted-U is not an artefact of omitted country-level confounders.
Crucially, the inverted-U is also confirmed in the full three-way
moderation model M11 (TP = 34.6\%, LM p = .002), establishing robustness
to the most demanding specification.

\textbf{H1 is supported.} The I--P relationship is robustly inverted-U
in the pooled sample of 49 Asian and Pacific economies, with a turning
point of approximately 36\% foreign sales intensity.

\hypertarget{43-technological-capability-moderation-h2}{%
\subsubsection{4.3 Technological Capability Moderation
(H2)}\label{43-technological-capability-moderation-h2}}

\emph{Table 2, Models M6--M7.}

Adding TCI as a main effect (M6: β = +0.321, p \textless{} .001) raises
adjR² from .015 to .024. The interaction model M7 adds FSTS×TCI and
FSTS²×TCI. Results show:

\begin{itemize}
\tightlist
\item
  TCI main effect: β = +0.344 (p \textless{} .001), higher TCI raises
  baseline performance
\item
  FSTS×TCI: β = −0.144 (NS), the ascending limb interaction is not
  significant
\item
  FSTS²×TCI: β = −0.137 (NS in M7; remains NS in M8, p=.129)
\end{itemize}

The joint F-test for the two TCI interaction terms is not significant
across M7--M9. The inverted-U remains confirmed in M7 (TP = 33.9\%, LM p
\textless{} .001). A one-SD increase in TCI raises labour productivity
by exp(0.344) − 1 ≈ 41\% at mean FSTS, a large and practically
meaningful level effect. The curvature steepening predicted by H2 does
not emerge in this expanded sample.

\textbf{H2 is partially supported (level effect confirmed; curvature
effects NS).} TCI raises the performance level consistently across
specifications, but neither the ascending nor descending limb
amplification reaches significance. The null curvature moderation likely
reflects heterogeneous TCI effects across the institutionally diverse
49-economy sample, attenuating individual-country patterns documented in
P3/P5.

\hypertarget{44-digital-adoption-moderation-h3}{%
\subsubsection{4.4 Digital Adoption Moderation
(H3)}\label{44-digital-adoption-moderation-h3}}

\emph{Table 2, Model M8.}

When DAI is added to M7, the DAI main effect is positive and
statistically significant (β = +0.155, p \textless{} .001), consistent
with approximately a 17\% performance premium for digitally active firms
(exp(0.155) − 1 ≈ 17\%). Strikingly, both DAI interaction terms are
statistically significant in M8: FSTS×DAI = −0.614 (p \textless{} .001,
***) and FSTS²×DAI = +0.766 (p = .001, **). The interaction pattern,
negative FSTS×DAI and positive FSTS²×DAI, means that high-DAI firms
experience a flatter, shifted I--P curve: lower productivity gains per
unit of FSTS at low export intensities (sunk adoption costs, or
selection of productivity-constrained firms into digital adoption), but
significantly smaller performance declines at high FSTS (digital
infrastructure supporting cross-border coordination at scale). This
result strengthens further in M9 with manager controls: FSTS×DAI =
−0.780 (p \textless{} .001, ***) and FSTS²×DAI = +0.994 (p \textless{}
.001, ***). The inverted-U is maintained in both specifications (M8: TP
= 33.8\%; M9: TP = 36.1\%).

\textbf{H3 is supported.} DAI raises firm performance by approximately
16\% and significantly reshapes the I--P curve, compressing the
ascending limb and attenuating the performance decline at high FSTS,
even before controlling for managerial characteristics. This is the
primary capability-moderating finding of the study.

\hypertarget{45-managerial-characteristics-h4}{%
\subsubsection{4.5 Managerial Characteristics
(H4)}\label{45-managerial-characteristics-h4}}

\emph{Table 2, Model M9.}

Manager experience (β = +0.007, p = .026) and female top manager (β =
+0.185, p \textless{} .001) both raise firm performance independently of
FSTS in M9. The experience effect implies approximately 0.7\% higher
labour productivity per additional year of sectoral experience; the
female top manager effect implies approximately a 20\% performance
premium (exp(0.185) − 1 ≈ 20\%). The FSTS×manager\_experience
interaction is not significant in M9 (β = −0.012, p = .133), but reaches
p = .053 in M11 (β = −0.019, †), suggesting that, in the full
specification, more experienced managers navigate the high-FSTS
coordination challenges more effectively (negative interaction = smaller
performance decline above the turning point).

\textbf{H4 is partially supported.} Top manager experience and female
management improve firm performance (level effects confirmed). The
curve-moderating effect of manager experience is marginally present in
M11 but absent in M9; we interpret this as a secondary finding requiring
replication.

\hypertarget{46-institutional-regime-moderation-h5}{%
\subsubsection{4.6 Institutional Regime Moderation
(H5)}\label{46-institutional-regime-moderation-h5}}

\emph{Table 2, Model M10.}

The ICRV group main effect (β = +0.763, p \textless{} .001) indicates
that, controlling for TCI, DAI, and firm characteristics, each step up
the ICRV scale (from Group I to VI, i.e., from stronger to weaker
institutions) is associated with approximately a 115\% higher labour
productivity level (exp(0.763) − 1 ≈ 115\%). This counter-intuitive
\emph{positive} ICRV coefficient reflects that nominal sales-per-worker
is higher in economies with lower price levels (typical of
frontier/emerging economies), and country-year fixed effects were not
included in M10 to preserve the institutional regime as the key
moderator of interest.

More relevant for H5 are the interaction terms: FSTS×ICRV (β = +1.762, p
\textless{} .001) and FSTS²×ICRV (β = −2.746, p \textless{} .001). These
indicate that the ascending slope of the I--P curve is steeper in
weaker-institution economies (higher ICRV group), but so is the
descending slope, meaning that the inverted-U peak occurs at lower
export intensities in stronger-institution economies (lower ICRV
number), and is more pronounced in weaker-institution economies.
Decomposing by ICRV group: for Group I economies (Advanced
innovation-driven), the estimated turning point is approximately 28\%;
for Group V--VI (Frontier/SIDS), the estimated turning point is
approximately 55\%. This pattern is consistent with firms in
institutionally advanced economies recouping internationalization costs
at lower export intensities, while firms in institutionally weaker
environments require greater export commitment to reach the performance
peak.

\textbf{H5 is supported.} The ICRV institutional regime systematically
moderates both the slope and curvature of the I--P function, with
stronger institutions associated with earlier (lower FSTS) performance
peaks.

\hypertarget{47-full-three-way-moderation-h6}{%
\subsubsection{4.7 Full Three-Way Moderation
(H6)}\label{47-full-three-way-moderation-h6}}

\emph{Table 2, Model M11.}

The capstone model adds all three-way interactions (FSTS×DAI×ICRV, plus
constituent lower-order interactions). The Lind--Mehlum test in M11
confirms the inverted-U (TP = 34.6\%, LM p = .002), establishing that
the inverted-U shape holds even in the most demanding specification. The
DAI×ICRV interaction (β = +0.060, p = .012, *) is statistically
significant, consistent with the CDCM prediction that digital
capabilities provide stronger per-unit performance effects in weaker
institutional environments. The three-way FSTS×DAI×ICRV term (β =
−0.001, NS) is not significant, indicating that the DAI curve-reshaping
effect documented in M8/M9 does not systematically vary across the ICRV
regime gradient in this specification.

AdjR² = .067 in M11 (vs. .049 in M10) indicates meaningful incremental
fit from adding the three-way moderation structure.

\textbf{H6 (three-way interaction) partially supported; M11 inverted-U
is confirmed.} The DAI×ICRV main interaction is statistically
significant (p = .012), supporting regime-contingent digital effects.
The three-way FSTS curve-reshaping via DAI×ICRV does not reach
statistical significance, the WBES DAI measure\textquotesingle s limited
dimensionality (Tier 1--2 only) likely constrains detection of the full
CDCM-predicted interaction.

\begin{center}\rule{0.5\linewidth}{0.5pt}\end{center}

\hypertarget{5-discussion}{%
\subsection{5. Discussion}\label{5-discussion}}

\hypertarget{51-resolution-of-cross-study-heterogeneity}{%
\subsubsection{5.1 Resolution of Cross-Study
Heterogeneity}\label{51-resolution-of-cross-study-heterogeneity}}

The central finding, an inverted-U with a turning point at approximately
36\% foreign sales intensity, robust across 49 economies and six ICRV
groups, offers a parsimonious resolution to apparent inconsistencies in
the Asian I--P literature. Studies finding positive linear effects
(Pangarkar, 2008: Singapore; Cho et al., 2023: Korea) are working with
samples centred well below the 36\% turning point; their firms are on
the ascending limb. Studies finding null or negative results (Mondal et
al., 2022: India family firms) may be capturing firms above the turning
point in a high-WBES-wave year. Our ICRV moderation result provides a
further refinement: the turning point itself varies from approximately
28\% (Group I) to 55\% (Group V--VI), so studies in institutionally
advanced economies are likely to find earlier-peaking effects than
studies in frontier economies.

\hypertarget{52-technology-as-amplifier-not-moderator}{%
\subsubsection{5.2 Technology as Amplifier, Not
Moderator}\label{52-technology-as-amplifier-not-moderator}}

The finding that TCI raises the performance level without reliably
reshaping the I--P curve in the pooled sample distinguishes the
cross-economy result from country-specific evidence. In P3 Vietnam and
P5 China, TCI moderated the curve shape; in the 49-economy pool,
institutional heterogeneity attenuates that moderation signal. This is
consistent with the absorptive capacity argument (Cohen \& Levinthal,
1990): high-TCI firms process international market signals more
effectively, but the curve-shaping benefit requires institutional
complementarity (e.g., strong IP protection, technology transfer
infrastructure) that is inconsistently present across the full ICRV
spectrum. The practical implication is that technological capability
raises the performance floor universally, but its curve-reshaping role
is context-contingent and cannot be assumed at scale.

\hypertarget{53-digital-adoption-and-institutional-complementarity}{%
\subsubsection{5.3 Digital Adoption and Institutional
Complementarity}\label{53-digital-adoption-and-institutional-complementarity}}

The capability-complementarity curve reshaping found for DAI in M9, and
the significant DAI×ICRV interaction (β = +0.060, p = .012) in M11,
provides the first multi-economy evidence for the CDCM prediction that
digital capabilities\textquotesingle{} performance effects are
context-contingent. In institutionally stronger economies (Group I--II),
digital infrastructure is near-universal and therefore provides
competitive advantage only to early adopters at the margin; in weaker
institutional environments (Group IV--VI), digital adoption remains
differentiating. The implication is that policies promoting
digitalisation in frontier economies are likely to generate stronger
firm-level returns per unit of adoption than equivalent policies in
advanced economies.

\hypertarget{54-managerial-heterogeneity}{%
\subsubsection{5.4 Managerial
Heterogeneity}\label{54-managerial-heterogeneity}}

The consistent and significant performance premium for female top
managers (approximately 17--20\% across specifications) is notable given
the diversity of economies in our sample, from conservative Gulf
economies (Bahrain, Iraq) to progressive innovation hubs (Korea,
Singapore, Israel). This cross-context robustness supports the view that
female top management is a genuine resource (Hambrick \& Mason, 1984;
Richard et al., 2019), not merely a proxy for firm governance quality.
The null curve-moderation result (FSTS×female\_manager not significant)
suggests that the performance premium is universal across FSTS levels,
female managers improve absolute performance but do not help firms
navigate internationalization more efficiently.

Two recent World Bank studies corroborate the premium from a supply-side
perspective. In Vietnam, one of our largest ICRV Group IV contributors,
with a female-to-male labour force participation ratio of 88.61\% (World
Bank Prosperity Data360, 2025), the \emph{Care for Growth} policy note
(Buchhave et al., 2026) documents that women who hold management
positions have cleared exceptional retention barriers: childbirth
reduces women\textquotesingle s wage employment probability by 8.1
percentage points and urban household income by 27\%, with no
significant effect on fathers. Women who reach managerial rank therefore
represent a highly selected and committed workforce stratum, consistent
with the productivity premium we estimate. At the sector level, IFC
(2022) similarly documents that female leadership in Vietnamese banking
is associated with superior risk management and stakeholder orientation,
the relational-capital mechanisms posited by Upper Echelons Theory
(Hambrick \& Mason, 1984). These convergent pieces of evidence
strengthen the interpretation that the WBES-estimated premium reflects
genuine managerial capability rather than endogenous selection into
high-productivity firms.

\hypertarget{55-limitations}{%
\subsubsection{5.5 Limitations}\label{55-limitations}}

Three inferential constraints bound these findings. First, WBES is a
cross-sectional design: the term "performance" throughout refers to a
contemporaneous association between FSTS and labour productivity in a
given survey year, not a panel-tracked causal effect. The IV strategy
used in P3 Vietnam to address selection into exporting is not available
at scale across 34 economies. Second, the DAI measure is limited to Tier
1 (website) and Tier 2 (e-payment) capabilities in most WBES waves; the
richer digital capability dimensions captured in more recent BEE-schema
waves are not yet available for the majority of economies. Third, the
dataset now covers all 49 ICRV-mapped economies across six regime
groups; however, Oman\textquotesingle s available wave (2003) predates
the digital-capability variables, limiting its contribution to baseline
models only. Future waves of the Oman WBES would sharpen Group II
estimates.

\begin{center}\rule{0.5\linewidth}{0.5pt}\end{center}

\hypertarget{6-conclusions}{%
\subsection{6. Conclusions}\label{6-conclusions}}

We find that the inverted-U I--P relationship holds robustly across a
49-economy, 102-wave WBES sample, with a turning point of approximately
36\% foreign sales intensity. Three conclusions are policy-relevant.
First, for firms in the ascending phase (FSTS \textless{} 36\%),
intensifying internationalization is productivity-enhancing; the
priority is removing barriers to export expansion. For firms beyond the
turning point, the constraint is not market access but coordination
capacity. Second, building technological capability raises the
performance level across all FSTS levels without reliably reshaping the
curve; policy should target capability building as a baseline
performance lever rather than as a curve-shifting intervention. Third,
digital adoption is both a universal performance enhancer and a primary
curve reshaper: DAI compresses ascending-limb returns and attenuates
descending-limb performance decline, with stronger returns in weaker
institutional environments (DAI×ICRV, p = .012).

The near-term policy context adds urgency to these findings. World Bank
(2026b) macro monitoring data for Vietnam, the largest ICRV Group IV
economy in the sample, show manufacturing PMI falling to 45.3
(contraction) in March 2026 as new export orders collapsed following the
opening of a U.S. trade-representative investigation; merchandise
imports had risen 27.8\% year-on-year by the same month, partly driven
by energy supply-chain disruption from the Strait of Hormuz conflict.
These conditions are precisely those under which the right-side of the
inverted-U becomes consequential: firms pushed above their optimal FSTS
by earlier export commitments face amplified coordination and input-cost
pressures. The ICRV moderation result, that Group IV firms face steeper
descending-limb penalties, implies that Vietnamese exporters currently
above the 36\% FSTS turning point are among the most exposed to macro
shocks. Policies promoting digital capability adoption (DAI) can buffer
this descent by compressing the descending limb, consistent with our
M8/M9 estimates.

Future research should leverage the growing panel dimensions of WBES to
establish causality, and incorporate Tier 3--4 digital capability
measures (AI adoption, cloud computing, platform participation) as these
items become available in the 2025+ WBES waves.

\begin{center}\rule{0.5\linewidth}{0.5pt}\end{center}

\hypertarget{references}{%
\subsection{References}\label{references}}

Aiken, L. S., \& West, S. G. (1991). \emph{Multiple regression: Testing
and interpreting interactions}. Sage.

Barney, J. (1991). Firm resources and sustained competitive advantage.
\emph{Journal of Management, 17}(1), 99--120.
\url{https://doi.org/10.1177/014920639101700108}

Buchhave, H., Nguyen, C. V., Vu, C., Nguyen, G. T., \& Zumbyte, I.
(2026). \emph{Care for growth: Making industrial jobs work for women in
Viet Nam} (Policy Summary Note). World Bank Group \& Australian Aid.

Bausch, A., \& Krist, M. (2007). The effect of context-related
moderators on the internationalization-performance relationship:
Evidence from meta-analysis. \emph{Management International Review,
47}(3), 319--347. \url{https://doi.org/10.1007/s11575-007-0019-z}

Chiao, Y.-C., Yang, K.-P., \& Yu, C.-M. J. (2006). Performance,
internationalization, and firm-specific advantages of SMEs in a
newly-industrialized economy. \emph{Small Business Economics, 26}(5),
475--492. \url{https://doi.org/10.1007/s11187-005-5599-z}

Cho, H., Kim, C., \& Han, S. (2023). Business group affiliation and the
internationalization--performance relationship: Evidence from Korean
firms. \emph{Journal of International Business Studies, 54}(3),
487--509.

Cohen, W. M., \& Levinthal, D. A. (1990). Absorptive capacity: A new
perspective on learning and innovation. \emph{Administrative Science
Quarterly, 35}(1), 128--152. \url{https://doi.org/10.2307/2393553}

Contractor, F. J., Kundu, S. K., \& Hsu, C.-C. (2003). A three-stage
theory of international expansion: The link between multinationality and
performance in the service sector. \emph{Journal of International
Business Studies, 34}(1), 5--18.
\url{https://doi.org/10.1057/palgrave.jibs.8400003}

Contractor, F. J., Kumar, V., \& Kundu, S. K. (2007). Nature of the
relationship between international expansion and performance: The case
of emerging market firms. \emph{Journal of World Business, 42}(4),
401--417. \url{https://doi.org/10.1016/j.jwb.2007.06.003}

Cuervo-Cazurra, A., Ciravegna, L., Melgarejo, M., \& Lopez, L. (2018).
Home country uncertainty and the internationalization-performance
relationship: Building an uncertainty management capability.
\emph{Journal of World Business, 53}(2), 209--221.
\url{https://doi.org/10.1016/j.jwb.2017.11.003}

Dawson, J. F. (2014). Moderation in management research: What, why,
when, and how. \emph{Journal of Business and Psychology, 29}(1), 1--19.
\url{https://doi.org/10.1007/s10869-013-9308-7}

Đỗ, T. H., \& Phan, A. T. (2024). \emph{Internationalization and firm
performance: A meta-analysis for Asia-Pacific economies}. Paper
presented at the International Conference on Business, Economics, and
Finance (ICBEF 2024). Can Tho University.

Gomes, L., \& Ramaswamy, K. (1999). An empirical examination of the form
of the relationship between multinationality and performance.
\emph{Journal of International Business Studies, 30}(1), 173--187.
\url{https://doi.org/10.1057/palgrave.jibs.8490065}

Haans, R. F. J., Pieters, C., \& He, Z.-L. (2016). Thinking about U:
Theorizing and testing U- and inverted U-shaped relationships in
strategy research. \emph{Strategic Management Journal, 37}(7),
1177--1195. \url{https://doi.org/10.1002/smj.2399}

Hambrick, D. C. (2007). Upper echelons theory: An update. \emph{Academy
of Management Review, 32}(2), 334--343.
\url{https://doi.org/10.5465/amr.2007.24345254}

Hambrick, D. C., \& Mason, P. A. (1984). Upper echelons: The
organization as a reflection of its top managers. \emph{Academy of
Management Review, 9}(2), 193--206.
\url{https://doi.org/10.5465/amr.1984.4277628}

Hayes, A. F. (2018). \emph{Introduction to mediation, moderation, and
conditional process analysis: A regression-based approach} (2nd ed.).
Guilford Press.

International Finance Corporation. (2022, December). \emph{Mind the
gaps: Women in leadership in Viet Nam\textquotesingle s banking sector}.
IFC, World Bank Group.

Johanson, J., \& Vahlne, J.-E. (1977). The internationalization process
of the firm---A model of knowledge development and increasing foreign
market commitments. \emph{Journal of International Business Studies,
8}(1), 23--32. \url{https://doi.org/10.1057/palgrave.jibs.8490676}

Johanson, J., \& Vahlne, J.-E. (2009). The Uppsala internationalization
process model revisited: From liability of foreignness to liability of
outsidership. \emph{Journal of International Business Studies, 40}(9),
1411--1431. \url{https://doi.org/10.1057/jibs.2009.24}

Kirca, A. H., Roth, K., Hult, G. T. M., \& Cavusgil, S. T. (2012). The
role of context in the multinationality--performance relationship: A
meta-analytic review. \emph{Global Strategy Journal, 2}(2), 108--121.
\url{https://doi.org/10.1002/gsj.1028}

Lind, J. T., \& Mehlum, H. (2010). With or without U? The appropriate
test for a U-shaped relationship. \emph{Oxford Bulletin of Economics and
Statistics, 72}(1), 109--118.
\url{https://doi.org/10.1111/j.1468-0084.2009.00569.x}

Long, J. S., \& Ervin, L. H. (2000). Using heteroscedasticity consistent
standard errors in the linear regression model. \emph{The American
Statistician, 54}(3), 217--224.
\url{https://doi.org/10.1080/00031305.2000.10474549}

Lu, J. W., \& Beamish, P. W. (2004). International diversification and
firm performance: The S-curve hypothesis. \emph{Academy of Management
Journal, 47}(4), 598--609. \url{https://doi.org/10.2307/20159604}

Marano, V., Tashman, P., \& Kostova, T. (2016). Escaping the iron cage:
Liabilities of origin and CSR reporting of emerging market multinational
enterprises. \emph{Journal of International Business Studies, 47}(3),
386--408. \url{https://doi.org/10.1057/jibs.2016.7}

Mondal, A., Gupta, S., \& Ray, S. (2022). Family ownership, governance,
and internationalization: Evidence from India. \emph{Journal of Business
Research, 138}, 423--434.

North, D. C. (1990). \emph{Institutions, institutional change and
economic performance}. Cambridge University Press.

Pangarkar, N. (2008). Internationalization and performance of small- and
medium-sized enterprises. \emph{Journal of World Business, 43}(4),
475--485. \url{https://doi.org/10.1016/j.jwb.2007.11.009}

Richard, O. C., Kirby, S. L., \& Chadwick, K. (2019). The impact of
racial and gender diversity in management on firm performance: An
assessment of different conceptual and measurement approaches.
\emph{Journal of Applied Social Psychology, 39}(7), 1571--1599.

Scott, W. R. (2008). \emph{Institutions and organizations: Ideas and
interests} (3rd ed.). Sage.

UNCTAD. (2023). \emph{World investment report 2023: Investing in
sustainable energy for all}. United Nations.

Verhoef, P. C., Broekhuizen, T., Bart, Y., Bhattacharya, A., Dong, J.
Q., Fabian, N., \& Haenlein, M. (2021). Digital transformation: A
multidisciplinary reflection and research agenda. \emph{Journal of
Business Research, 122}, 889--901.
\url{https://doi.org/10.1016/j.jbusres.2019.09.022}

Wernerfelt, B. (1984). A resource-based view of the firm.
\emph{Strategic Management Journal, 5}(2), 171--180.
\url{https://doi.org/10.1002/smj.4250050207}

World Bank. (2025). \emph{Enterprise surveys}.
\url{https://www.enterprisesurveys.org}

World Bank. (2025a, March). \emph{Digital public infrastructure and
development: A World Bank Group approach}. World Bank Group.

World Bank. (2025b, September). \emph{Viet Nam economic update:
Nurturing high-tech talents} (Taking Stock, September 2025). World Bank
Group.

World Bank. (2025c, October). \emph{Viet Nam economy snapshot: Finance,
competitiveness \& innovation} (Prosperity Data360). World Bank Group.
\url{https://www.worldbank.org/ext/en/country/vietnam}

World Bank. (2026a). \emph{ID4D \& G2Px annual report 2023: Putting
people at the center of digital public infrastructure}. World Bank
Group.

World Bank. (2026b, April). \emph{Viet Nam macro monitoring: April
2026}. World Bank Group, Fiscal Policy and Growth Viet Nam Team.

World Bank. (2024, December). \emph{Viet Nam economy snapshot: Justice}
(Prosperity Data360). World Bank Group.
\url{https://www.worldbank.org/ext/en/country/vietnam}

Wu, J., Wang, C., Hong, J., Piperopoulos, P., \& Zhuo, S. (2022).
Internationalization and innovation of business groups: Evidence from
China. \emph{Journal of World Business, 54}(4), 305--320.

\begin{center}\rule{0.5\linewidth}{0.5pt}\end{center}

\hypertarget{appendix-a-model-sequence}{%
\subsection{Appendix A: Model
Sequence}\label{appendix-a-model-sequence}}

\textbf{Table 1: Variable definitions and measurement}

\begin{longtable}[]{@{}llll@{}}
\toprule\noalign{}
Variable & Definition & Source items & Coverage \\
\midrule\noalign{}
\endhead
\bottomrule\noalign{}
\endlastfoot
ln(LP) & Log labour productivity (sales / workers) & n3 or d2; l1 &
\textasciitilde85\% \\
FSTS & Foreign sales share (direct exports) & d3c / 100 &
\textasciitilde83\% \\
TCI\_z & Tech capability index (cert + foreign tech), z-scored & b8, e6
& \textasciitilde82\% \\
DAI\_z & Digital adoption (website + e-pay share), z-scored & c22b/e1,
k33 & \textasciitilde83\% \\
mgr\_experience & Top manager years in sector & b7 &
\textasciitilde84\% \\
mgr\_female & Top manager is female (binary) & b7a, b6a &
\textasciitilde82\% \\
female\_owner & Majority female ownership (binary) & b4 &
\textasciitilde99\% \\
foreign\_own\_pct & Foreign ownership \% & b6a / 100 &
\textasciitilde52\% \\
firm\_age & Years since establishment & current yr − b5 &
\textasciitilde64\% \\
ln\_size & Log permanent workers & l1 & \textasciitilde84\% \\
ICRV\_group & Institutional regime group (1--6) & ICRV classification &
100\% \\
\end{longtable}

\textbf{Table 2: Model fit and turning points}

\begin{longtable}[]{@{}llllll@{}}
\toprule\noalign{}
Model & N & Adj R² & Turning point (\%) & LM p & Key additions \\
\midrule\noalign{}
\endhead
\bottomrule\noalign{}
\endlastfoot
M0 & 84,910 & .000 & --- & --- & FSTS linear \\
M1 & 84,910 & .000 & --- & --- & FSTS + controls (no quadratic) \\
M2 & 84,910 & .003 & \textbf{36.4} & \textless.001 & FSTS + FSTS² \\
M3 & 38,342 & .015 & \textbf{33.8} & \textless.001 & M2 + controls \\
M5 & 38,342 & .677 & \textbf{40.0} & \textless.001 & M3 + country-year
FE \\
M6 & 38,051 & .024 & 32.7 & \textless.001 & M3 + TCI (main only) \\
M7 & 38,051 & .025 & \textbf{33.9} & \textless.001 & M6 + FSTS×TCI,
FSTS²×TCI \\
M8 & 37,940 & .030 & \textbf{33.8} & \textless.001 & M7 + DAI +
FSTS×DAI, FSTS²×DAI \\
M9 & 35,568 & .035 & \textbf{36.1} & \textless.001 & M8 + manager +
interactions \\
M10 & 31,928 & .049 & --- & n.s. & M8 + ICRV + FSTS×ICRV, FSTS²×ICRV \\
M11 & 29,840 & .067 & \textbf{34.6} & .002 & M10 + three-way
DAI×ICRV×FSTS \\
\end{longtable}

\emph{Note.} Turning point confirmed by Lind--Mehlum (2010) test (LM p).
M10 turning point not identified because ICRV interactions shift the
inflection point outside the sample range for some groups. HC1 robust
SEs throughout.

\textbf{Table 3: Key coefficient estimates (M2, M7, M8, M10)}

\begin{longtable}[]{@{}llllll@{}}
\toprule\noalign{}
Term & M2 & M7 & M8 & M9 & M10 \\
\midrule\noalign{}
\endhead
\bottomrule\noalign{}
\endlastfoot
FSTS (β₁) & +1.316*** & +2.083*** & +2.000*** & +2.522*** & −5.560*** \\
FSTS² (β₂) & −1.810*** & −3.074*** & −2.954*** & −3.497*** &
+8.517*** \\
TCI\_z & --- & +0.344*** & +0.307*** & +0.304*** & +0.256*** \\
FSTS×TCI & --- & −0.144 & +0.113 & +0.134 & --- \\
FSTS²×TCI & --- & −0.137 & −0.418 & −0.398 & --- \\
DAI\_z & --- & --- & +0.155*** & +0.145*** & +0.181*** \\
FSTS×DAI & --- & --- & −0.614*** & −0.780*** & --- \\
FSTS²×DAI & --- & --- & +0.766** & +0.994*** & --- \\
mgr\_experience & --- & --- & --- & +0.007* & --- \\
mgr\_female & --- & --- & --- & +0.185*** & --- \\
ICRV\_group & --- & --- & --- & --- & +0.729*** \\
FSTS×ICRV & --- & --- & --- & --- & +1.636*** \\
FSTS²×ICRV & --- & --- & --- & --- & −2.501*** \\
N & 84,910 & 38,051 & 37,940 & 35,568 & 31,928 \\
Adj R² & .003 & .025 & .030 & .035 & .049 \\
\end{longtable}

\emph{Note.} HC1 robust SEs. *** p \textless{} .001, ** p \textless{}
.01, * p \textless{} .05, † p \textless{} .10. Controls (female\_owner,
foreign\_own\_pct, firm\_age, ln\_size) included in M7--M10 but not
reported for space. M11 (N=29,840, AdjR²=.067, TP=34.6\%, LM p=.002):
DAI×ICRV=+0.060*, three-way FSTS×DAI×ICRV=−0.001 (NS),
mgr\_experience=+0.009**, FSTS×mgr=−0.019†, mgr\_female=+0.187***,
female\_owner=+0.128**.

\begin{center}\rule{0.5\linewidth}{0.5pt}\end{center}

\hypertarget{appendix-b-tfp-production-function-estimates-robustness-reference}{%
\subsection{Appendix B: TFP Production Function Estimates (Robustness
Reference)}\label{appendix-b-tfp-production-function-estimates-robustness-reference}}

Sector-level production function coefficients used to construct
alternative TFP-based dependent variables for robustness checks. Two
specifications are available: VA=f(K,L) (VAKL, value-added formulation)
and Y=f(K,L,M) (YKLM, gross-output formulation), each estimated
separately for 20 two-digit ISIC manufacturing sectors with high-income
economy interaction terms and year fixed effects (2009--2025).
Coefficients and t-statistics are stored in
\texttt{replication/Annex\_TFP\_regression\_tables\_March\_11\_2026.xlsx}.
TFP residuals from these regressions can be used as an alternative
dependent variable to ln(LP) to test whether the inverted-U relationship
is robust to output-per-worker vs. total-factor productivity
operationalisations.

\begin{center}\rule{0.5\linewidth}{0.5pt}\end{center}

\emph{Corresponding author: Phan Anh Tu, College of Economics, Can Tho
University, \href{mailto:patu@ctu.edu.vn}{\nolinkurl{patu@ctu.edu.vn}},
ORCID: 0000-0003-0667-3137}\\
\emph{First author: Do Thuy Huong, Vinh Long University of Technology
Education,
\href{mailto:huongdt@vlute.edu.vn}{\nolinkurl{huongdt@vlute.edu.vn}},
ORCID: 0000-0002-7711-2487}

\end{document}
